\documentclass[aspectratio=169, 11pt, xcolor={table}]{beamer}
% Shared style for the course slides, based on presentation_text/main.tex.
% Dark blue removed: every slide (cover included) uses the light background
% of the template's content slides, with teal + orange as accents.

\usepackage{amsmath}
\usepackage{amssymb}
\usepackage{booktabs}
\usepackage{colortbl}
\usepackage{tabularx}
\usepackage{tikz}
\usepackage[skins]{tcolorbox}
\usepackage{fontawesome5}
\usepackage{polyglossia}

\usetikzlibrary{positioning, calc, arrows.meta, shapes.geometric}

%------------------------------------------------
% Languages and fonts (compile from slides/ with XeLaTeX or tectonic)
%------------------------------------------------
\setdefaultlanguage[numerals=maghrib]{arabic}
\setotherlanguage{english}
\usefonttheme{serif}

\newfontfamily\arabicfont[
    Script=Arabic, Scale=0.98, Path=../fonts/, Extension=.ttf,
    UprightFont=*-Regular, BoldFont=*-Bold,
    ItalicFont=*-Italic, SlantedFont=*-Italic,
    BoldItalicFont=*-BoldItalic, BoldSlantedFont=*-BoldItalic]{Amiri}
\newfontfamily\arabicfontsf[
    Script=Arabic, Scale=0.98, Path=../fonts/, Extension=.ttf,
    UprightFont=*-Regular, BoldFont=*-Bold,
    ItalicFont=*-Italic, SlantedFont=*-Italic,
    BoldItalicFont=*-BoldItalic, BoldSlantedFont=*-BoldItalic]{Amiri}
% Heading font
\newfontfamily\khatt[Script=Arabic, Scale=1.05, Path=../fonts/,
    BoldFont=ae_AlBattar.ttf]{ae_AlBattar.ttf}

%------------------------------------------------
% Colors (template palette without the navy tones)
%------------------------------------------------
\definecolor{brandTeal}{HTML}{00B4D8}
\definecolor{brandOrange}{HTML}{FF7043}
\definecolor{brandLight}{HTML}{F8F9FA}
\definecolor{brandDark}{HTML}{263238}
\definecolor{brandMuted}{HTML}{6C757D}
\colorlet{tealDeep}{brandTeal!70!black}

\setbeamercolor{background canvas}{bg=brandLight}
\setbeamercolor{normal text}{fg=brandDark}
\setbeamercolor{structure}{fg=brandTeal}
\setbeamertemplate{navigation symbols}{}

% beamer.cls sets \raggedright globally; under RTL that centers short lines
% or pushes them left, so go back to justified text (flush right).
\rightskip=0pt\relax
% Lists call \raggedright too; make it a no-op for RTL slides.
\AtBeginDocument{\renewcommand{\raggedright}{\rightskip=0pt\relax}}

%------------------------------------------------
% Frame title, footer, lists
%------------------------------------------------
\setbeamertemplate{frametitle}{%
    \vspace{0.35cm}
    \noindent{\khatt\bfseries\Large\color{brandDark}\insertframetitle}
    \par
    \vskip-2pt
    \noindent{\color{brandTeal}\rule{\textwidth}{1.5pt}}
    \vspace{0.05cm}
}

% Footer: page number only
\setbeamertemplate{footline}{%
    \hbox to \paperwidth{\hfill
        {\color{brandMuted}\scriptsize\LR{\insertframenumber\,/\,\inserttotalframenumber}}%
        \hspace{0.6cm}}%
    \vspace{0.25cm}}

\setbeamertemplate{itemize item}{\color{brandTeal}\small\faChevronLeft}
\setbeamertemplate{itemize subitem}{\color{brandMuted}\small\faAngleLeft}
\setbeamertemplate{enumerate item}{\color{tealDeep}\bfseries\insertenumlabel)}

% Tables: every cell framed by a visible border
\newcolumntype{R}{>{\raggedleft\arraybackslash}X}
\newcolumntype{C}{>{\centering\arraybackslash}X}
\newcommand{\ras}[1]{\cellcolor{brandTeal!20}\textbf{#1}}
% Colored text inside a table cell (\leavevmode keeps it on the row baseline)
\newcommand{\cell}[2]{\leavevmode\textcolor{#1}{#2}}
\renewcommand{\arraystretch}{1.35}
\setlength{\arrayrulewidth}{0.6pt}
\arrayrulecolor{brandMuted!70}

%------------------------------------------------
% Boxes (the two boxes of the template's 2nd slide)
%------------------------------------------------
\tcbset{
    brandbox/.style={
        enhanced, arc=4pt, outer arc=4pt, boxrule=1.5pt,
        coltitle=white, coltext=brandDark, fonttitle=\khatt\bfseries,
        toptitle=0.8mm, bottomtitle=0.8mm,
        top=1.2mm, bottom=1.2mm, left=3mm, right=3mm,
        before skip=4pt, after skip=4pt}
}
% Teal box: definitions, objectives, summaries
\newtcolorbox{tbox}[2][]{brandbox, colback=brandTeal!10,
    colframe=brandTeal, title={#2}, #1}
% Orange box: warnings
\newtcolorbox{obox}[2][]{brandbox, colback=brandOrange!10,
    colframe=brandOrange, title={#2}, #1}
% Card (template's 3rd slide, in teal)
\newtcolorbox{card}[2][]{enhanced, colback=brandTeal!6, colframe=brandTeal,
    arc=6pt, boxrule=1.2pt, halign=center, halign title=center,
    coltitle=white, coltext=brandDark, fonttitle=\khatt\bfseries,
    top=1.5mm, bottom=1.5mm, left=2mm, right=2mm, title={#2}, #1}
% Prompt text written to the tool
\newtcolorbox{amr}[1][brandTeal]{enhanced, arc=1mm, boxrule=0pt,
    colback=white, overlay={\fill[#1] (frame.north east) rectangle ([xshift=-2.5pt]frame.south east);},
    left=6pt, right=8pt, top=2pt, bottom=2pt,
    before skip=3pt, after skip=3pt}

%------------------------------------------------
% TikZ (Arabic right-to-left inside nodes)
%------------------------------------------------
\newcommand{\nodecolor}{}
\tikzset{
    every node/.append style={
        execute at begin node={\nodecolor\everypar{\beginR}\beginR},
        execute at end node=\endR},
    lawn/.code={\def\nodecolor{\color{#1}}},
    khatwa/.style={draw=#1, fill=#1!12, very thick, rounded corners=5pt,
        minimum height=0.95cm, text width=1.9cm, align=center,
        font=\khatt\bfseries\small},
    khatwa/.default=brandTeal,
    sahm/.style={-{Stealth[length=2.5mm]}, thick, color=brandMuted},
    rang/.style={draw=#1, fill=#1!12},
    raqm/.style={circle, fill=#1, lawn=white, inner sep=1.5pt,
        font=\small\bfseries},
}

% Horizontal chain of steps, right to left:
% \flow[width]{{a}, {b\\c}, {d}}
\newcommand{\flow}[2][1.9cm]{%
    \begin{tikzpicture}[node distance=5mm]
        \foreach \t [count=\i] in {#2} {
            \ifnum\i=1
                \node[khatwa, text width=#1] (f\i) {\t};
            \else
                \pgfmathtruncatemacro\p{\i-1}
                \node[khatwa, text width=#1, left=of f\p] (f\i) {\t};
                \draw[sahm] (f\p) -- (f\i);
            \fi
        }
    \end{tikzpicture}}

%------------------------------------------------
% Cover (template's 1st slide on the light background)
%------------------------------------------------
\newcommand{\cover}[3]{%
\begin{frame}[plain]
    \begin{tikzpicture}[remember picture, overlay]
        \draw[brandTeal, line width=1.5pt, opacity=0.3]
            (current page.south east) ++(-1,1) circle (4cm);
        \draw[brandTeal, line width=0.8pt, opacity=0.25]
            (current page.north west) ++(2,2) circle (6cm);
        \draw[brandOrange, line width=0.5pt, opacity=0.25]
            (current page.north west) ++(2,2) circle (7cm);
        \node[anchor=north east] at ($(current page.north east)+(-1.0,-0.8)$) {
            \begin{tikzpicture}[scale=0.35]
                \fill[brandTeal] (0,0) -- (0.8,1.2) -- (1.6,0) -- cycle;
                \fill[brandOrange] (0.8,1.2) -- (1.6,2.4) -- (2.4,1.2) -- cycle;
            \end{tikzpicture}
        };
    \end{tikzpicture}

    \vspace{-0.05cm}
    {\hspace*{1.1cm}\khatt\bfseries\color{brandDark}
        المدرسة العليا للأساتذة ببوزريعة ــ ملحقة بومرداس\par}

    \vspace{1.2cm}
    \begin{minipage}{0.85\textwidth}
        {\color{tealDeep}\khatt\bfseries #1} \\[0.3cm]
        {\raggedleft\color{brandDark}\Huge\khatt\bfseries #2\par} \vspace{0.3cm}
        {\color{brandMuted}\large #3}
    \end{minipage}

    \vspace{1.2cm}
    \begin{minipage}{0.85\textwidth}
        \color{brandMuted}\small
        \begin{tabular}{@{}ll}
            \textbf{الأستاذ:} & د. بن عبد الرزاق يوسف \\
            \textbf{الفئة:}   & طلبة قسم اللغة العربية \\
        \end{tabular}
    \end{minipage}
\end{frame}}

\title{تصميم مساعد ذكاء اصطناعي وسير عمل}

\begin{document}

\cover{درس الفصل السابع}
{تصميم مساعد ذكاء اصطناعي وسير عمل}
{تعليمات ــ معرفة ــ اختبار ــ تحسين ــ سير عمل}

\begin{frame}{أهداف الدرس}
    \begin{tbox}{\faBullseye\quad أهداف الدرس}
        في نهاية هذا الدرس، تستطيع أن:
        \begin{itemize}
            \item تشرح ما المساعد المخصص، وما مكوناته.
            \item تكتب تعليمات واضحة لمساعد تعليمي.
            \item تختبر المساعد وتحسّن تعليماته.
            \item تقسّم مهمة كبيرة إلى سير عمل من عدة مراحل.
        \end{itemize}
    \end{tbox}
\end{frame}

\begin{frame}{لماذا نبني مساعدًا مخصصًا؟}
    يستعمل الأستاذ الذكاء الاصطناعي كل أسبوع، لكنه يكتب التعليمات نفسها في كل محادثة،
    ويرفع الملفات نفسها في كل مرة.
    الحل: نبني \textbf{مساعدًا مخصصًا} مرة واحدة، ثم نستعمله كلما احتجنا إليه.
    \begin{tbox}{\faLightbulb\quad تعريف}
        \textbf{المساعد المخصص:} نسخة من أداة الذكاء الاصطناعي،
        نعطيها اسمًا وتعليمات ثابتة وملفات مرجعية، لتقوم بمهمة محددة.
    \end{tbox}
    نبنيه في \textenglish{ChatGPT} باستعمال \textenglish{GPTs}،
    أو في \textenglish{Claude} باستعمال \textenglish{Projects}.
\end{frame}

\begin{frame}{مكونات المساعد}
    \begin{center}
    \begin{tikzpicture}[
        juz/.style={draw=brandTeal, fill=brandTeal!10, thick, rounded corners=5pt,
            text width=4.3cm, align=center, minimum height=1.35cm}]
        \node[circle, draw=brandOrange, fill=brandOrange!15, very thick,
            text width=1.7cm, align=center, font=\khatt\bfseries] (m) {المساعد};
        \node[juz] (a) at (4.1,1.1)
            {{\khatt\bfseries الاسم}\\\small يعرّف المساعد ومهمته من أول نظرة};
        \node[juz] (b) at (4.1,-1.1)
            {{\khatt\bfseries التعليمات}\\\small تحدد دوره وطريقة عمله وحدوده};
        \node[juz] (c) at (-4.1,1.1)
            {{\khatt\bfseries المعرفة}\\\small ملفات مرجعية: دروس، ملخصات، محاضرات};
        \node[juz] (d) at (-4.1,-1.1)
            {{\khatt\bfseries الأدوات}\\\small قدرات إضافية، نفعّلها عند الحاجة فقط};
        \foreach \x in {a,b,c,d} \draw[thick, brandMuted] (m) -- (\x);
    \end{tikzpicture}
    \end{center}
\end{frame}

\begin{frame}{كتابة التعليمات}
    \small
    التعليمات هي «عقل» المساعد. نكتبها بعناصر الأمر الجيد (الفصل الثالث):
    \begin{amr}[brandTeal]
        \textbf{\color{tealDeep}الدور:} أنت مساعد تعليمي متخصص في مادة النحو العربي.\\
        \textbf{\color{tealDeep}المهمة:} ساعد الطلبة على فهم الدروس بطريقة بسيطة ومنظمة.\\
        \textbf{\color{tealDeep}طريقة العمل:} اشرح خطوة بخطوة، واستعمل أمثلة بسيطة.\\
        \textbf{\color{tealDeep}المراجع:} اعتمد على الملفات المرفقة عندما يتعلق السؤال بالمادة.\\
        \textbf{\color{brandOrange}القيود:} لا تخترع معلومات. إذا لم تجد المعلومة، صرّح بذلك.\\
        \textbf{\color{tealDeep}شكل الإجابة:} عناوين واضحة، ونقاط قصيرة، وأمثلة.\\
        \textbf{\color{tealDeep}السؤال الغامض:} اطلب توضيحًا قبل أن تجيب.
    \end{amr}
\end{frame}

\begin{frame}{دورة بناء المساعد}
    بناء المساعد لا ينتهي بإنشائه: نختبره، ونكتشف أخطاءه، ونحسّنه، ثم نختبره من جديد.
    \begin{center}
    \begin{tikzpicture}[
        marhala/.style={thick, rounded corners=4pt,
            text width=2.1cm, align=center, minimum height=0.8cm,
            font=\khatt\bfseries\small}]
        \node[marhala, rang=brandTeal]   (a) at (90:2.3)   {حدّد المهمة};
        \node[marhala, rang=brandTeal]   (b) at (30:2.8)   {اكتب التعليمات};
        \node[marhala, rang=brandTeal]   (c) at (-30:2.8)  {أنشئ وأضف الملفات};
        \node[marhala, rang=brandTeal]   (d) at (-90:2.3)  {اختبر};
        \node[marhala, rang=brandOrange] (e) at (-150:2.8) {اكتشف الأخطاء};
        \node[marhala, rang=brandTeal]   (f) at (150:2.8)  {حسّن التعليمات};
        \draw[sahm] (a) to[bend left=20] (b);
        \draw[sahm] (b) to[bend left=20] (c);
        \draw[sahm] (c) to[bend left=20] (d);
        \draw[sahm] (d) to[bend left=20] (e);
        \draw[sahm] (e) to[bend left=20] (f);
        \draw[sahm, dashed] (f) to[bend left=10]
            node[above, sloped, font=\scriptsize, lawn=brandMuted] {أعد الاختبار} (d);
    \end{tikzpicture}
    \end{center}
\end{frame}

\begin{frame}{اختبار المساعد}
    لا نفترض أن المساعد يعمل جيدًا. نختبره بأربعة أنواع من الأسئلة:
    \begin{center}
    \small
    \begin{tabularx}{\textwidth}{|>{\centering\arraybackslash}p{3cm}|R|R|}
        \hline
        \ras{الاختبار} & \ras{مثال} & \ras{ماذا نلاحظ؟} \\ \hline
        سؤال بسيط             & ما المبتدأ؟ اشرحه لطالب مبتدئ. & هل اللغة بسيطة؟ \\ \hline
        سؤال يعتمد على الملف   & اعتمد على الملف، واذكر أهم ثلاثة مفاهيم. & هل الإجابة مطابقة للملف؟ \\ \hline
        سؤال ناقص             & اشرح لي هذا الموضوع بالتفصيل. & هل يطلب توضيحًا؟ \\ \hline
        سؤال خارج المهمة       & كيف أحضّر البيتزا؟ & هل يحافظ على دوره؟ \\ \hline
    \end{tabularx}
    \end{center}
\end{frame}

\begin{frame}{تحسين المساعد بعد الاختبار}
    لاحظنا أن المساعد يعطي إجابات طويلة جدًا. نضيف إلى التعليمات:
    \begin{amr}[brandTeal]
        استعمل جملًا قصيرة ونقاطًا واضحة.
        لا تقدم شرحًا طويلًا إلا إذا طلب المستخدم ذلك.
    \end{amr}
    ثم نعيد الاختبار نفسه، ونقارن الإجابة \textbf{قبل التعديل وبعده}.
\end{frame}

\begin{frame}{سير العمل \textenglish{(Workflow)}}
    \begin{tbox}{\faLightbulb\quad تعريف}
        \textbf{سير العمل:} تقسيم مهمة كبيرة إلى مراحل صغيرة متتابعة.
        مخرج كل مرحلة يصبح مدخلًا للمرحلة التالية.
    \end{tbox}
    \textbf{\color{tealDeep}مثال: من درس طويل إلى أسئلة تدريبية}
    \begin{center}
        \flow[1.9cm]{{النص\\الأصلي}, {استخراج\\الأفكار}, {إنشاء\\الملخص}, {إنشاء\\الأسئلة}, {التحقق\\النهائي}}
    \end{center}
    \begin{center}
        \cell{brandOrange}{\faUserCheck}\quad \textbf{مراجعة بشرية بعد كل مرحلة}
    \end{center}
\end{frame}

\begin{frame}{مراحل سير العمل}
    \small
    \begin{center}
    \begin{tabularx}{\textwidth}{|>{\centering\arraybackslash}p{1.6cm}|C|C|C|}
        \hline
        \ras{المرحلة} & \ras{المدخل} & \ras{المعالجة} & \ras{المخرج} \\ \hline
        1 & نص المحاضرة & استخراج الأفكار الأساسية & قائمة أفكار \\ \hline
        2 & قائمة الأفكار & كتابة ملخص منظم & ملخص الدرس \\ \hline
        3 & الملخص & إنشاء خمسة أسئلة مع الإجابات & أسئلة تدريبية \\ \hline
    \end{tabularx}
    \end{center}
    \begin{obox}{\faExclamationTriangle\quad انتبه!}
        لا ننتقل إلى المرحلة التالية قبل مراجعة نتيجة المرحلة الحالية.
        الخطأ الصغير في البداية يكبر في المراحل التالية.
    \end{obox}
\end{frame}

\begin{frame}{أمر واحد أم سير عمل؟}
    \begin{center}
    \begin{tabularx}{0.85\textwidth}{|R|C|C|}
        \hline
        \ras{المعيار} & \ras{أمر واحد} & \ras{سير عمل} \\ \hline
        سهولة المراجعة         & \cell{brandOrange}{صعبة}   & \cell{tealDeep}{\textbf{سهلة}} \\ \hline
        اكتشاف الأخطاء         & \cell{brandOrange}{صعب}    & \cell{tealDeep}{\textbf{سهل}} \\ \hline
        التحكم في كل مرحلة     & \cell{brandOrange}{ضعيف}   & \cell{tealDeep}{\textbf{قوي}} \\ \hline
        إعادة الاستعمال        & \cell{brandMuted}{محدودة}  & \cell{tealDeep}{\textbf{ممكنة}} \\ \hline
        السرعة                 & \cell{tealDeep}{\textbf{أسرع}} & \cell{brandMuted}{أبطأ قليلًا} \\ \hline
    \end{tabularx}
    \end{center}
\end{frame}

\begin{frame}{أفكار لمساعدين في قسم اللغة العربية}
    \begin{tbox}{\faChalkboardTeacher\quad أفكار للتطبيق}
        \begin{itemize}
            \item \textbf{مساعد الإعراب:} يشرح إعراب الجمل خطوة بخطوة، اعتمادًا على كتاب النحو المقرر.
            \item \textbf{مساعد البلاغة:} يستخرج الصور البيانية من نص، ويشرحها.
            \item \textbf{مساعد التعبير:} يقترح ملاحظات على إنشاء التلميذ دون أن يكتبه بدلًا منه.
        \end{itemize}
    \end{tbox}
\end{frame}

\begin{frame}{خلاصة الدرس}
    \begin{tbox}{\faCheckCircle\quad خلاصة الدرس}
        \begin{itemize}
            \item المساعد المخصص = \textbf{اسم} + \textbf{تعليمات} + \textbf{معرفة} + \textbf{أدوات}.
            \item نبني، ثم \textbf{نختبر}، ثم \textbf{نحسّن}، ثم نختبر من جديد.
            \item \textbf{سير العمل} يقسم المهمة الكبيرة إلى مراحل، مع مراجعة بعد كل مرحلة.
            \item المساعد الجيد يحتاج إلى أستاذ يراجع نتائجه باستمرار.
        \end{itemize}
    \end{tbox}
\end{frame}

\end{document}
